\documentclass[11pt]{article}
\usepackage[margin=1in]{geometry}
\usepackage{enumitem}
% =============================================================================
% Preambles/header.tex — shared LaTeX/Beamer preamble for this project
%
% Use from a Beamer lecture by prepending:
%     \input{header}
% and compiling with TEXINPUTS=../Preambles:$TEXINPUTS (the /compile-latex
% skill does this automatically).
%
% PALETTE CONTRACT. Color names defined here MUST match the names in
% Quarto/theme-template.scss so Beamer and Quarto renderings use the same
% palette. The scripts/check-palette-sync.sh script enforces this.
%
% To customize: edit the HEX values below AND the matching $variable in
% Quarto/theme-template.scss, then run scripts/check-palette-sync.sh.
% =============================================================================

% -----------------------------------------------------------------------------
% Engine detection + encoding
%
% Our /compile-latex skill uses XeLaTeX, where inputenc is unnecessary and
% can produce encoding warnings. Only load inputenc under pdfTeX.
% -----------------------------------------------------------------------------
\usepackage{iftex}
\ifPDFTeX
  \usepackage[utf8]{inputenc}
\fi

\usepackage{amsmath, amssymb, amsthm}
\usepackage{booktabs}
\usepackage{graphicx}

% -----------------------------------------------------------------------------
% PALETTE — mirrors Quarto/theme-template.scss variable names.
% Load xcolor and define colors BEFORE anything that references them
% (notably hyperref, hypersetup, and the Beamer color assignments).
% -----------------------------------------------------------------------------
\usepackage{xcolor}

\definecolor{primary-blue}{HTML}{012169}      % headings, accents
\definecolor{primary-gold}{HTML}{B9975B}      % emphasis, borders
\definecolor{highlight-yellow}{HTML}{F2A900}  % markers, alerts
\definecolor{light-bg}{HTML}{E8EDF5}          % title / section backgrounds
\definecolor{jet}{HTML}{1A1A1A}               % body text

% Semantic colors (used in diagrams and annotations)
\definecolor{positive}{HTML}{15803D}          % good / observed / identified
\definecolor{negative}{HTML}{B91C1C}          % bad / problematic / bias
\definecolor{neutral}{HTML}{525252}           % reference / context

% Highlight accents (match .hi-* classes in the SCSS)
\definecolor{hi-slate}{HTML}{314F4F}
\definecolor{hi-green}{HTML}{2E8B57}
\definecolor{hi-red}{HTML}{C41E3A}

% Semantic aliases so skill templates and snippets can write
% \textcolor{accent}{...} without hard-coding "primary-blue".
\colorlet{accent}{primary-blue}
\colorlet{accent2}{primary-gold}

% -----------------------------------------------------------------------------
% Hyperref — loaded AFTER xcolor + palette so link colors resolve.
% -----------------------------------------------------------------------------
\usepackage{hyperref}
\hypersetup{colorlinks=true, linkcolor=primary-blue, urlcolor=primary-blue,
            citecolor=primary-blue, pdfborder={0 0 0}}

% -----------------------------------------------------------------------------
% Course code — set once per deck (after \input{header}, before \begin{document})
% via \coursecode{CS401}. Shown in the footer on every slide so a deck's course
% is identifiable without opening the title page. Empty by default.
% -----------------------------------------------------------------------------
\makeatletter
\newcommand{\coursecode}[1]{\gdef\@coursecodetext{#1}}
\gdef\@coursecodetext{}
\makeatother

% -----------------------------------------------------------------------------
% Beamer theme assignments (only apply under Beamer).
%
% We detect Beamer via \@ifclassloaded{beamer}, which is the documented,
% non-internal way. This must be wrapped in \makeatletter/\makeatother
% because \@ifclassloaded contains an @-catcode character.
% -----------------------------------------------------------------------------
\makeatletter
\@ifclassloaded{beamer}{%
  \usefonttheme{professionalfonts}
  \setbeamercolor{structure}{fg=primary-blue}
  \setbeamercolor{frametitle}{fg=primary-blue}
  \setbeamercolor{title}{fg=primary-blue}
  \setbeamercolor{subtitle}{fg=primary-gold}
  \setbeamercolor{author}{fg=primary-blue}
  \setbeamercolor{institute}{fg=neutral}
  \setbeamercolor{date}{fg=primary-gold}
  \setbeamercolor{itemize item}{fg=primary-blue}
  \setbeamercolor{itemize subitem}{fg=primary-gold}
  \setbeamercolor{alerted text}{fg=primary-gold}
  \setbeamercolor{block title}{fg=white, bg=primary-blue}
  \setbeamercolor{block body}{bg=light-bg}
  % Semantic block variants: block=definition/concept (blue), exampleblock=
  % worked example (green/positive), alertblock=key takeaway or caution (gold).
  % Keep to <=2 colored boxes per slide across ALL three variants (INV-7).
  \setbeamercolor{block title example}{fg=white, bg=positive}
  \setbeamercolor{block body example}{bg=light-bg}
  \setbeamercolor{block title alerted}{fg=white, bg=primary-gold}
  \setbeamercolor{block body alerted}{bg=light-bg}

  % Minimal footer: course code (left) + page X / N (right), no navigation clutter
  \setbeamertemplate{navigation symbols}{}
  \setbeamertemplate{footline}{%
    \color{neutral}\footnotesize\hspace{0.7em}\@coursecodetext\hfill
    \insertframenumber/\inserttotalframenumber\hspace{0.7em}\vspace{0.3em}}
}{}
\makeatother

% -----------------------------------------------------------------------------
% TikZ setup — libraries the snippet gallery uses
% -----------------------------------------------------------------------------
\usepackage{tikz}
\usetikzlibrary{arrows.meta, positioning, calc, decorations.pathreplacing,
                fit, shapes.geometric, backgrounds}

% Shared TikZ styles — snippets and hand-written diagrams can pick these up
% via `\begin{tikzpicture}[every node/.style={...}]` or by naming specific
% styles (e.g., `\node[dag-node] (X) at (0,0) {X};`).
\tikzset{
  % Node styles for causal DAGs and similar
  dag-node/.style = {
    draw=primary-blue, fill=light-bg, rounded corners=2pt,
    minimum width=1.2cm, minimum height=0.9cm, align=center, font=\small
  },
  decision-node/.style = {
    draw=primary-gold, fill=white, diamond, aspect=1.8,
    minimum width=1.8cm, minimum height=1.2cm, align=center, font=\small,
    inner sep=1pt
  },
  % Edge styles — observed vs counterfactual
  observed-edge/.style = {-{Latex[length=2.5mm]}, thick, draw=jet},
  counterfactual-edge/.style = {-{Latex[length=2.5mm]}, thick, draw=neutral, dashed},
  confound-edge/.style = {-{Latex[length=2.5mm]}, thick, draw=negative, dashed},
  % Data-point styles for plots
  observed-dot/.style = {fill=primary-blue, draw=primary-blue, circle, inner sep=1.2pt},
  counterfactual-dot/.style = {fill=white, draw=neutral, circle, inner sep=1.2pt}
}

% -----------------------------------------------------------------------------
% Convenience macros
% -----------------------------------------------------------------------------
\newcommand{\muted}[1]{\textcolor{neutral}{#1}}
\newcommand{\key}[1]{\textcolor{primary-gold}{\textbf{#1}}}
\newcommand{\good}[1]{\textcolor{positive}{#1}}
\newcommand{\bad}[1]{\textcolor{negative}{#1}}

% Standout frame macro: full-bleed dark background for section transitions.
% Beamer-only (uses \begin{frame}/\setbeamercolor/\usebeamerfont) -- do not
% call from an article-class file (e.g. Notes/<CODE>/*-notes.tex). \muted,
% \key, \good, \bad, and \coursecode above are plain xcolor/gdef and are
% safe to use from either class; verified when Notes/article-class support
% was added.
% Expand: \transitionslide{Your transition title}
\newcommand{\transitionslide}[1]{%
  \begingroup
  \setbeamercolor{background canvas}{bg=primary-blue}
  \begin{frame}[plain]
    \centering
    \vfill
    {\usebeamerfont{title}\color{white}\Large #1\par}
    \vfill
  \end{frame}
  \endgroup
}

% Section divider frame (Beamer-only), an alternative to \transitionslide with a
% darker two-line style: near-black (jet) background, a small white label line
% above a larger gold title, and a thin gold rule along the bottom edge.
% Expand: \sectiondivider{Section 2}{Pointers}
\newcommand{\sectiondivider}[2]{%
  \begingroup
  \setbeamercolor{background canvas}{bg=jet}
  \begin{frame}[plain]
    \vspace*{3.0cm}
    {\color{white}\ttfamily\large #1\par}
    \vspace{0.5cm}
    {\color{primary-gold}\LARGE #2\par}
    \begin{tikzpicture}[remember picture, overlay]
      \draw[primary-gold, line width=2.5pt]
        ($(current page.south west)+(0,0.1)$) -- ($(current page.south east)+(0,0.1)$);
    \end{tikzpicture}
  \end{frame}
  \endgroup
}

% -----------------------------------------------------------------------------
% Channel watermark — "Concepts That Click" (YouTube).
%
% Article-class files (CompetitiveExam Q/A sets, Notes, Assignments) get a
% light diagonal draft watermark on every page plus a centered footer naming
% the channel. Beamer decks get a subtle TikZ background watermark instead
% (the footline already carries the course code). Centralized here so every
% MCQ PDF is branded without per-file edits.
% -----------------------------------------------------------------------------
\makeatletter
\@ifclassloaded{article}{%
  \usepackage{draftwatermark}
  \SetWatermarkText{Concepts That Click}
  \SetWatermarkScale{1}
  \SetWatermarkFontSize{2cm}
  \SetWatermarkLightness{0.86}
  \SetWatermarkAngle{45}
  \usepackage{fancyhdr}
  \pagestyle{fancy}
  \fancyhf{}
  \fancyfoot[C]{\footnotesize\textcolor{neutral}{Concepts That Click~|~YouTube: \href{https://www.youtube.com/@concepts.thatclick}{@concepts.thatclick}}}
  \renewcommand{\headrulewidth}{0pt}
  \renewcommand{\footrulewidth}{0pt}
  \fancypagestyle{plain}{%
    \fancyhf{}%
    \fancyfoot[C]{\footnotesize\textcolor{neutral}{Concepts That Click~|~YouTube: \href{https://www.youtube.com/@concepts.thatclick}{@concepts.thatclick}}}%
    \renewcommand{\headrulewidth}{0pt}%
    \renewcommand{\footrulewidth}{0pt}%
  }
}{}
\@ifclassloaded{beamer}{%
  \setbeamertemplate{background}{%
    \begin{tikzpicture}[remember picture, overlay]
      \node[rotate=45, opacity=0.055, align=center]
        at (current page.center)
        {\Huge\textcolor{neutral}{Concepts That Click}};
    \end{tikzpicture}%
  }
}{}
\makeatother 
\coursecode{JKSSB}
\renewcommand{\thesection}{16.\arabic{section}}
\newtheorem{example}{Example}[section]
\newenvironment{definitionbox}[1]{%
  \par\noindent\textbf{Definition (#1).}\ \itshape
}{\par\vspace{0.3em}}
\title{Computer Ports --- Detailed Lecture Notes}
\author{JKSSB Computer Awareness}
\date{\today}

\begin{document}
\maketitle

\section{What a port is}
Every device that connects to a computer does so through a \textbf{port}. A
port is a socket or interface on the machine into which a cable or a device
is plugged. It lets data and, in many cases, electric power pass between the
computer and the attached device. Ports are not all alike: they differ in
their shape, in the type of signal they carry (digital or analog), and in the
job they are meant to do.

The common physical ports found on a personal computer are USB (including
USB-C), HDMI, DisplayPort, VGA, Ethernet through an RJ-45 connector, the
3.5~mm audio jack, and SATA. Knowing which port suits which device is the
central skill of this lesson.

\begin{definitionbox}{Port}
A socket or interface on a device where a cable or another device is
connected, allowing data (and often power) to flow between the computer and
the attached device.
\end{definitionbox}

\begin{center}
\begin{tabular}{p{0.26\textwidth}p{0.64\textwidth}}
\toprule
\textbf{Term} & \textbf{Meaning in this lesson} \\
\midrule
Port & The socket on the device (a USB port, an HDMI port). \\
Connector & The plug end of a cable that fits the port. \\
Cable & The wire that joins two connectors. \\
Physical port & A hardware socket on the machine. \\
Logical port & A software number that names a network service. \\
\bottomrule
\end{tabular}
\end{center}

A first view of the ports covered in this lesson, with their signal type and
their normal use, is given below. The rest of the lesson explains each one in
turn.

\begin{center}
\begin{tabular}{p{0.20\textwidth}p{0.22\textwidth}p{0.48\textwidth}}
\toprule
\textbf{Port} & \textbf{Signal} & \textbf{Typical use} \\
\midrule
USB / USB-C & Digital (data + power) & Keyboard, mouse, pen drive, phone \\
HDMI & Digital (video + audio) & TV, monitor, projector \\
DisplayPort & Digital (video + audio) & PC monitors, graphics cards \\
VGA & Analog (video only) & Older monitors and projectors \\
RJ-45 (Ethernet) & Digital (network) & Wired LAN cable, router, switch \\
3.5 mm audio jack & Analog (audio) & Headphones, speakers, microphone \\
SATA / eSATA & Digital (storage) & Internal and external drives \\
\bottomrule
\end{tabular}
\end{center}

\section{The two meanings of ``port''}
The word \textbf{port} has two entirely different meanings in computing, and
the exam frequently tests the difference. A \textbf{physical port} is a piece
of hardware --- a socket such as a USB port, an HDMI port, or an RJ-45
network port. A \textbf{logical port}, also called a network port, is a
software number that identifies a particular service running on a computer
during network communication.

A web server, for example, listens on logical port 80 for ordinary web
requests and on port 443 for secure web requests. The physical network cable,
meanwhile, is plugged into an RJ-45 port on the network card. Thus one
machine uses both meanings at the same time without any contradiction: the
socket is hardware, and the port number is software.

\begin{definitionbox}{Logical (network) port}
A number, in the range 0 to 65535, that identifies a specific service or
process on a computer for the purpose of network communication. It is not a
physical socket.
\end{definitionbox}

\begin{example}
An item states, ``The monitor is connected to logical port 80.'' This is
nonsense. Port 80 is a logical port used by web (HTTP) traffic, not a socket
for a monitor. A monitor connects to a physical video port such as HDMI,
DisplayPort, or VGA. The wrong statement has mixed the two meanings of the
word \emph{port}.
\end{example}

A few well-known logical port numbers are worth remembering, because the exam
sometimes asks which service uses a given port. The list below is
illustrative; the concept being tested is that these are \emph{numbers}, not
sockets.

\begin{center}
\begin{tabular}{p{0.30\textwidth}p{0.16\textwidth}p{0.40\textwidth}}
\toprule
\textbf{Service} & \textbf{Logical port} & \textbf{Full form} \\
\midrule
HTTP & 80 & HyperText Transfer Protocol \\
HTTPS & 443 & HTTP Secure \\
FTP & 21 & File Transfer Protocol \\
SMTP & 25 & Simple Mail Transfer Protocol \\
DNS & 53 & Domain Name System \\
\bottomrule
\end{tabular}
\end{center}

\section{USB: the Universal Serial Bus}
\textbf{USB} stands for \textbf{Universal Serial Bus}. It is a standard
interface for connecting a wide range of peripherals to a computer. Two of
its best-known features are \textbf{plug and play} and \textbf{hot swapping}:
a device can be attached or removed while the computer is running, without
switching the machine off. A USB connection carries both \textbf{data} and
\textbf{power}, which is why many small devices can be used or charged
directly from a USB port.

Typical USB devices include the keyboard, the mouse, a printer, a pen drive
(flash drive), an external hard disk, a webcam, and a smartphone. The
connector at the end of a USB cable comes in several shapes:

\begin{itemize}
  \item \textbf{USB Type-A}: the familiar rectangular connector found on most
    PCs.
  \item \textbf{USB Type-B}: often used for printers and scanners.
  \item \textbf{Mini-USB} and \textbf{Micro-USB}: smaller connectors used on
    older phones and devices.
  \item \textbf{USB Type-C}: a modern, small, reversible connector.
\end{itemize}

\textbf{USB-C}, or USB Type-C, deserves special attention. It is a
\emph{connector type}, not a separate bus. It is \textbf{reversible}, so it
can be plugged in either way up, and a single USB-C connection can carry
data, power, and video together. It is now common on smartphones, laptops,
external SSDs, and docking stations. The key exam point is precision: USB is
the standard, USB-C is one connector shape for that standard, and a USB port
is the socket on the machine.

The USB standard itself has developed over time through versions such as USB
1.1, USB 2.0, USB 3.0 (and its later 3.1 and 3.2 refinements), and USB4. Each
newer version generally offers a higher data-transfer rate, and the version
is independent of the connector shape: a Type-A or Type-C socket can appear
on a machine that supports different USB versions. For exam purposes, the
safe points are the meaning of the abbreviation, the plug-and-play and
hot-swap behaviour, the combined data-and-power carrying, and the
reversibility of USB-C.

\section{Video ports: HDMI, DisplayPort, and VGA}
Displays are attached through video ports. Three names dominate the syllabus.

\textbf{HDMI} stands for \textbf{High-Definition Multimedia Interface}. It
carries \textbf{digital video and digital audio} over a single cable, which
is why one HDMI cable can replace a separate video cable and a separate audio
cable. It is very common on televisions, monitors, projectors, set-top boxes,
and game consoles.

\textbf{DisplayPort} is also a \textbf{digital} audio/video interface, but it
is designed mainly for computers and monitors. It appears on graphics cards,
laptops, and desktop monitors, and it supports high refresh rates and
multiple monitors through daisy-chaining. \textbf{Mini DisplayPort} is a
smaller variant of the same interface.

\begin{center}
\begin{tabular}{p{0.22\textwidth}p{0.30\textwidth}p{0.38\textwidth}}
\toprule
\textbf{Point} & \textbf{HDMI} & \textbf{DisplayPort} \\
\midrule
Signal & Digital video + audio & Digital video + audio \\
Main use & TV, home theatre, monitor & PC monitors, graphics cards \\
Note & Very common on TVs & Common on PCs; multi-monitor \\
\bottomrule
\end{tabular}
\end{center}

Both HDMI and DisplayPort are digital. The right choice in practice usually
depends on which sockets the devices offer, not on one being inherently
better than the other.

\textbf{VGA} stands for \textbf{Video Graphics Array}. It is an
\textbf{analog} video interface: it carries video only and no audio. It uses
a 15-pin D-sub connector, often coloured blue and held in place by two
screws. VGA is still seen on older monitors and projectors but has largely
been replaced by HDMI and DisplayPort.

\begin{definitionbox}{Digital vs analog video}
Digital video ports (HDMI, DisplayPort) send the picture as discrete numbers
and also carry audio; the analog port (VGA) sends a continuously varying
signal and carries video only.
\end{definitionbox}

\begin{example}
An item reads, ``VGA carries both video and audio.'' This is false. VGA is an
analog video-only interface with no audio channel. HDMI and DisplayPort, by
contrast, carry digital video together with digital audio. The wrong option
has transferred the audio capability of the digital ports onto VGA.
\end{example}

\section{Ethernet and the RJ-45 connector}
\textbf{Ethernet} is the standard for \textbf{wired} computer networking, and
it is commonly associated with the IEEE 802.3 family of standards. The
physical socket used for a wired network connection is the \textbf{RJ-45}
connector. RJ-45 stands for \textbf{Registered Jack 45}, and it has eight
pins, used with twisted-pair cable such as Cat5e or Cat6.

Typical devices attached through an RJ-45 port are a router, a switch, a
modem, and the network card of a desktop computer. A near-neighbour that
causes confusion is \textbf{RJ-11}, the smaller telephone jack. RJ-45 is the
network connector with eight pins; RJ-11 is the telephone connector and is
smaller. The two look similar but serve different purposes.

\section{Audio jacks}
The common audio socket on a personal computer is the \textbf{3.5~mm audio
jack}. It connects headphones, earphones, speakers, and microphones. The plug
wiring has standard forms: \textbf{TRS} (Tip-Ring-Sleeve) carries stereo
audio, while \textbf{TRRS} (Tip-Ring-Ring-Sleeve) adds a microphone channel,
which is why a phone headset with a microphone uses a TRRS plug. On desktop
PCs a colour convention is often followed: green for audio output, such as
headphones or speakers, and pink for a microphone. Some equipment uses
2.5~mm or 6.35~mm jacks instead of the 3.5~mm size.

\section{Storage interfaces: SATA and eSATA}
\textbf{SATA} stands for \textbf{Serial Advanced Technology Attachment}. It
is the interface that connects \textbf{internal storage} --- a hard disk
drive, a solid-state drive, or an optical drive --- to the motherboard. The
SATA data cable is thin and has seven pins, and a separate power connector
feeds the drive. \textbf{eSATA} stands for \textbf{external SATA}, and it is
used to connect an external drive. It is important to classify eSATA
correctly: it is a \emph{storage interface}, not a general-purpose peripheral
port like USB. SATA is about attaching drives, whereas USB, HDMI, and the
audio jack are about attaching other kinds of peripherals and displays.

\section{Classifying a port: four questions to ask}
When a question shows a port or names it, four questions settle almost every
case. First, \emph{is it a physical port or a logical one?} A socket is
physical; a number such as 80 or 443 is logical. Second, \emph{what kind of
signal does it carry?} HDMI and DisplayPort are digital audio/video, VGA is
analog video, the 3.5~mm jack is analog audio, and RJ-45 is a digital network
link. Third, \emph{what class of device does it attach?} Displays go to
video ports, drives go to SATA, network equipment goes to RJ-45, and general
peripherals go to USB. Fourth, \emph{which near-neighbour is it often
confused with?} USB-C is confused with USB generally; RJ-45 is confused with
RJ-11; VGA is confused with HDMI and DisplayPort.

\begin{definitionbox}{Port class}
A way of grouping ports by what they attach: peripheral ports (USB), display
ports (HDMI, DisplayPort, VGA), the network port (RJ-45), audio ports (3.5~mm
jack), and storage interfaces (SATA, eSATA).
\end{definitionbox}

Two ports are easily confused with a near-neighbour. The table below keeps the
distinctions in one place.

\begin{center}
\begin{tabular}{p{0.16\textwidth}p{0.26\textwidth}p{0.44\textwidth}}
\toprule
\textbf{Port} & \textbf{Nearest confusion} & \textbf{The distinction} \\
\midrule
USB-C & USB (the standard) & USB-C is a connector shape; USB is the bus. \\
HDMI & DisplayPort & Both digital A/V; HDMI is common on TVs, DisplayPort on PCs. \\
VGA & HDMI / DisplayPort & VGA is analog, video-only; the others are digital with audio. \\
RJ-45 & RJ-11 & RJ-45 is the 8-pin network jack; RJ-11 is the smaller phone jack. \\
SATA & USB & SATA attaches drives (storage); USB is a general peripheral port. \\
Logical port & Physical port & Logical port is a number; physical port is a socket. \\
\bottomrule
\end{tabular}
\end{center}

\section{Exam retrieval and common confusions}
The following distinctions prevent most errors on this topic.
\begin{enumerate}
  \item A \textbf{physical port} is a hardware socket; a \textbf{logical
    (network) port} is a software number, such as HTTP 80 or HTTPS 443. Never
    mix the two meanings of ``port''.
  \item \textbf{VGA is analog} and video-only; \textbf{HDMI and DisplayPort
    are digital} and carry both video and audio.
  \item \textbf{RJ-45} (network, eight pins) is not \textbf{RJ-11} (the
    smaller telephone jack).
  \item \textbf{USB-C is a connector type} and is reversible; ``USB'' is the
    standard, and a USB port is the socket.
  \item \textbf{SATA is a storage interface}; do not call it a display or
    network port.
  \item Say the term precisely: a \emph{port} is the socket, a
    \emph{connector} is the plug, and a \emph{cable} is the wire. This is the
    same discipline as file \emph{format} vs file \emph{extension} (TRM-09).
\end{enumerate}

\begin{example}
An item asks which port carries digital audio and digital video over one
cable. Name the digital audio/video ports: HDMI and DisplayPort both qualify,
whereas VGA is analog video-only, RJ-45 is networking, and SATA is storage.
If the stem prefers the port common on television sets, choose HDMI; if it
stresses PC monitors and multi-monitor setups, choose DisplayPort.
\end{example}

The full forms below are the ones most likely to be asked directly, so they
should be memorised exactly.

\begin{center}
\begin{tabular}{p{0.18\textwidth}p{0.56\textwidth}}
\toprule
\textbf{Abbreviation} & \textbf{Full form} \\
\midrule
USB & Universal Serial Bus \\
HDMI & High-Definition Multimedia Interface \\
VGA & Video Graphics Array \\
RJ-45 & Registered Jack 45 \\
SATA & Serial Advanced Technology Attachment \\
eSATA & external SATA \\
TRS & Tip-Ring-Sleeve \\
TRRS & Tip-Ring-Ring-Sleeve \\
HTTP & HyperText Transfer Protocol \\
HTTPS & HTTP Secure \\
\bottomrule
\end{tabular}
\end{center}

\subsection*{Exercises}
\begin{enumerate}[label=\arabic*.]
  \item Distinguish a physical port from a logical port, giving one example
    of each.
  \item Expand USB, HDMI, VGA, RJ-45, and SATA.
  \item Which video port is analog and video-only, and which two common video
    ports are digital and also carry audio?
  \item What is the difference between RJ-45 and RJ-11?
  \item What is the difference between SATA and eSATA, and why is SATA not a
    general peripheral port?
  \item Which port is commonly used to connect a monitor and supports high
    refresh rates and multiple monitors, and which connector is used for a
    wired Ethernet connection?
  \item Explain the difference between TRS and TRRS audio plugs.
  \item Name two features of USB that make attaching devices convenient.
\end{enumerate}

\subsection*{Solutions}
\begingroup\small
\begin{enumerate}[label=\arabic*.]
  \item A physical port is a hardware socket such as a USB port or an RJ-45
    port. A logical port is a software number that identifies a network
    service, such as HTTP port 80.
  \item USB = Universal Serial Bus; HDMI = High-Definition Multimedia
    Interface; VGA = Video Graphics Array; RJ-45 = Registered Jack 45;
    SATA = Serial Advanced Technology Attachment.
  \item VGA is the analog, video-only port. HDMI and DisplayPort are digital
    and carry both video and audio.
  \item RJ-45 is the eight-pin connector used for wired Ethernet networking;
    RJ-11 is the smaller telephone jack.
  \item SATA connects internal drives (HDD, SSD, optical drive) to the
    motherboard. eSATA is external SATA, used for an external drive. SATA is a
    storage interface --- it attaches drives --- so it is not a
    general-purpose peripheral port such as USB.
  \item DisplayPort is the monitor port for high refresh rates and multiple
    monitors; RJ-45 is the connector for a wired Ethernet connection.
  \item TRS (Tip-Ring-Sleeve) has three contact areas and carries stereo
    audio; TRRS (Tip-Ring-Ring-Sleeve) has four and adds a microphone channel.
  \item USB supports plug and play and hot swapping, and it carries both data
    and power over the same connection.
\end{enumerate}
\endgroup
\end{document}
